\documentclass[11pt,a4paper]{article}

\usepackage[utf8]{inputenc}
\usepackage[T1]{fontenc}
\usepackage[french]{babel}
\usepackage[margin=2.3cm]{geometry}
\usepackage{booktabs}
\usepackage{pgfplots}
\usepackage{pgfplotstable}
\pgfplotsset{compat=1.18}

\title{Comparaison d'un arbre rouge-noir avec \texttt{std::map}}
\author{TP de C++ avancé}
\date{\today}

\newif\ifdonneesMSI
% Passer a \donneesMSItrue apres avoir rempli resultats_benchmark.csv.
\donneesMSIfalse

\pgfplotstableread[col sep=comma]{resultats_benchmark.csv}\benchmarkdata

\begin{document}
\maketitle

\begin{abstract}
Ce TP compare les temps d'insertion, de recherche et de suppression d'un
arbre rouge-noir C++ avec ceux de \texttt{std::map<int,int>}. Les mesures
sont prises sur plusieurs tailles d'entree avec des cles aleatoires.
\end{abstract}

\section{Environnement}
\begin{itemize}
  \item Ordinateur : MSI Stealth GS66 12UGS.
  \item Processeur : Intel Core i7-12700H, 14 coeurs (6 P-cores et 8 E-cores),
        20 threads.
  \item Memoire : 16 Go DDR5.
  \item Carte graphique : NVIDIA GeForce RTX 3070 Ti, non utilisee.
  \item Compilation : MinGW-w64 / g++, C++17, optimisation \texttt{-O2}.
        Version de g++ : \emph{a relever sur la machine de mesure}.
\end{itemize}
Le benchmark est mono-thread CPU. Les resultats doivent etre releves sur le
MSI ; aucun temps provenant d'une autre machine n'est reporte ici.

\section{Protocole}
Le programme genere des cles aleatoires distinctes et emploie les memes cles
et valeurs dans les deux structures. Pour chaque taille
$N \in \{1\,000,\ 10\,000,\ 100\,000,\ 1\,000\,000\}$, il mesure :
\begin{itemize}
  \item l'insertion avec \texttt{emplace} ;
  \item la recherche de toutes les cles avec \texttt{find} ;
  \item la suppression de toutes les cles avec \texttt{erase(cle)}.
\end{itemize}
Chaque operation est repetee cinq fois apres une chauffe. L'ordre des deux
structures alterne et le resultat affiche est la mediane en nanosecondes par
operation. La mesure utilise \texttt{std::chrono::steady\_clock}. Les
invariants de l'arbre sont verifies hors chronometrage.

Les medianes affichees sont a reporter dans \texttt{resultats\_benchmark.csv}.
Les graphiques restent volontairement vides tant que ces mesures ne sont pas
renseignees. Apres avoir rempli le CSV, remplacer
\texttt{\textbackslash donneesMSIfalse} par
\texttt{\textbackslash donneesMSItrue} ci-dessus, puis compiler le rapport.

\section{Resultats}
\begin{figure}[ht]
\centering
\begin{tikzpicture}
\begin{axis}[
  width=0.9\textwidth,height=6cm,
  xmode=log,log basis x=10,
  xmin=1000,xmax=1000000,
  xtick={1000,10000,100000,1000000},
  xlabel={$N$},ylabel={ns par operation},
  title={Insertion},grid=major,
  scaled x ticks=false,
]
\ifdonneesMSI
  \addplot+[mark=*] table[x=N,y=RBTreeInsertion]{\benchmarkdata};
  \addlegendentry{RBTree}
  \addplot+[mark=square*] table[x=N,y=std_mapInsertion]{\benchmarkdata};
  \addlegendentry{\texttt{std::map}}
\else
  \node at (rel axis cs:0.5,0.5) {Medianes MSI a renseigner dans le CSV};
\fi
\end{axis}
\end{tikzpicture}

\vspace{0.6cm}
\begin{tikzpicture}
\begin{axis}[
  width=0.9\textwidth,height=6cm,
  xmode=log,log basis x=10,
  xmin=1000,xmax=1000000,
  xtick={1000,10000,100000,1000000},
  xlabel={$N$},ylabel={ns par operation},
  title={Recherche},grid=major,
  scaled x ticks=false,
]
\ifdonneesMSI
  \addplot+[mark=*] table[x=N,y=RBTreeFind]{\benchmarkdata};
  \addlegendentry{RBTree}
  \addplot+[mark=square*] table[x=N,y=std_mapFind]{\benchmarkdata};
  \addlegendentry{\texttt{std::map}}
\else
  \node at (rel axis cs:0.5,0.5) {Medianes MSI a renseigner dans le CSV};
\fi
\end{axis}
\end{tikzpicture}

\vspace{0.6cm}
\begin{tikzpicture}
\begin{axis}[
  width=0.9\textwidth,height=6cm,
  xmode=log,log basis x=10,
  xmin=1000,xmax=1000000,
  xtick={1000,10000,100000,1000000},
  xlabel={$N$},ylabel={ns par operation},
  title={Suppression},grid=major,
  scaled x ticks=false,
]
\ifdonneesMSI
  \addplot+[mark=*] table[x=N,y=RBTreeErase]{\benchmarkdata};
  \addlegendentry{RBTree}
  \addplot+[mark=square*] table[x=N,y=std_mapErase]{\benchmarkdata};
  \addlegendentry{\texttt{std::map}}
\else
  \node at (rel axis cs:0.5,0.5) {Medianes MSI a renseigner dans le CSV};
\fi
\end{axis}
\end{tikzpicture}
\caption{Temps median par operation en fonction du nombre de cles.}
\end{figure}

\section{Conclusion}
Les deux structures equilibrent leur arbre et les operations mesurees ont une
complexite logarithmique en moyenne. La comparaison experimentale depend des
constantes d'implementation, du compilateur et du processeur. Aucune structure
ne peut donc etre declaree plus rapide sur ce MSI avant d'avoir complete les
mesures et les graphiques. Le rapport \texttt{std::map}/RBTree est superieur
a 1 lorsque \texttt{std::map} est plus lent, et inferieur a 1 lorsqu'il est
plus rapide.

\end{document}
